% from ega1-5-fr.tex
\oldpage[I]{138}

\begin{proposition}[5.5.5]
\label{I.5.5.5-fr}
Soit $(Y_\lambda)$ un recouvrement ouvert d'un préschéma $Y$~; pour qu'un
morphisme $f:X\to Y$ soit séparé, il faut et il suffit que chacune de ses
restrictions $f^{-1}(Y_\lambda)\to Y_\lambda$ soit séparée.
\end{proposition}

Si on pose $X_\lambda=f^{-1}(Y_\lambda)$, tout revient, compte tenu de
(\hyperref[I.4.2.4-fr]{4.2.4, b)}) et de l'identité des produits
$X_\lambda\times_Y X_\lambda$ et
$X_\lambda\times_{Y_\lambda}X_\lambda$
(\hyperref[I.3.2.5-fr]{3.2.5}), à prouver que les
$X_\lambda\times_Y X_\lambda$ forment un recouvrement de $X\times_Y X$. Or,
si l'on pose $Y_{\lambda\mu}=Y_\lambda\cap Y_\mu$ et
$X_{\lambda\mu}=X_\lambda\cap X_\mu=f^{-1}(Y_{\lambda\mu})$,
$X_\lambda\times_Y X_\mu$ s'identifie au produit
$X_{\lambda\mu}\times_{Y_{\lambda\mu}}X_{\lambda\mu}$
(\hyperref[I.3.2.6.4-fr]{3.2.6.4}), donc aussi à
$X_{\lambda\mu}\times_Y X_{\lambda\mu}$
(\hyperref[I.3.2.5-fr]{3.2.5}), et finalement à un ouvert de
$X_\lambda\times_Y X_\lambda$, ce qui établit notre assertion
(\hyperref[I.3.2.7-fr]{3.2.7}).

La prop. (\hyperref[I.5.5.4-fr]{5.5.4}) permet, en prenant un recouvrement de
$Y$ par des ouverts affines, de ramener l'étude des morphismes séparés à celle
des morphismes séparés à valeurs dans des schémas affines.

\begin{proposition}[5.5.6]
\label{I.5.5.6-fr}
Soient $Y$ un schéma affine, $X$ un préschéma, $(U_\alpha)$ un recouvrement de
$X$ par des ouverts affines. Pour qu'un morphisme $f:X\to Y$ soit séparé, il
faut et il suffit que, pour tout couple d'indices $(\alpha,\beta)$,
$U_\alpha\cap U_\beta$ soit un ouvert affine, et que l'anneau
$\Gamma(U_\alpha\cap U_\beta,\mathscr{O}_X)$ soit engendré par la réunion des
images canoniques des anneaux $\Gamma(U_\alpha,\mathscr{O}_X)$ et
$\Gamma(U_\beta,\mathscr{O}_X)$.
\end{proposition}

Les $U_\alpha\times_Y U_\beta$ forment un recouvrement ouvert de
$X\times_Y X$ (\hyperref[I.3.2.7-fr]{3.2.7})~; en désignant par $p$ et $q$ les
projections de $X\times_Y X$, on a
\[
  \begin{aligned}
    \Delta_X^{-1}(U_\alpha\times_Y U_\beta)
      &=\Delta_X^{-1}\bigl(p^{-1}(U_\alpha)\cap q^{-1}(U_\beta)\bigr)\\
      &=\Delta_X^{-1}\bigl(p^{-1}(U_\alpha)\bigr)
        \cap\Delta_X^{-1}\bigl(q^{-1}(U_\beta)\bigr)
        =U_\alpha\cap U_\beta~;
  \end{aligned}
\]
tout revient donc à exprimer que la restriction de $\Delta_X$ à
$U_\alpha\cap U_\beta$ est une immersion fermée dans
$U_\alpha\times_Y U_\beta$. Or, cette restriction n'est autre que
$(j_\alpha,j_\beta)_Y$, en désignant par $j_\alpha$ (resp. $j_\beta$) le
morphisme d'injection de $U_\alpha\cap U_\beta$ dans $U_\alpha$ (resp.
$U_\beta$), ainsi qu'il résulte des définitions. Comme
$U_\alpha\times_Y U_\beta$ est un schéma affine dont l'anneau est
canoniquement isomorphe à
$\Gamma(U_\alpha,\mathscr{O}_X)
 \otimes_{\Gamma(Y,\mathscr{O}_Y)}
 \Gamma(U_\beta,\mathscr{O}_X)$
(\hyperref[I.3.2.2-fr]{3.2.2}), on voit que $U_\alpha\cap U_\beta$ doit être
un schéma affine et que l'application
$h_\alpha\otimes h_\beta\to h_\alpha h_\beta$ de l'anneau
$A(U_\alpha\times_Y U_\beta)$ dans
$\Gamma(U_\alpha\cap U_\beta,\mathscr{O}_X)$ doit être surjective
(\hyperref[I.4.2.3-fr]{4.2.3}), ce qui achève la démonstration.

\begin{corollary}[5.5.7]
\label{I.5.5.7-fr}
Un schéma affine est séparé (et est par suite un schéma, ce qui justifie la
terminologie de (\hyperref[I.5.4.1-fr]{5.4.1})).
\end{corollary}

\begin{corollary}[5.5.8]
\label{I.5.5.8-fr}
Soit $Y$ un schéma affine~; pour que $f:X\to Y$ soit un morphisme séparé, il
faut et il suffit que $X$ soit séparé (autrement dit, que $X$ soit un schéma).
\end{corollary}

On observera en effet que le critère de
(\hyperref[I.5.5.6-fr]{5.5.6}) ne dépend pas de $f$.

\begin{corollary}[5.5.9]
\label{I.5.5.9-fr}
Pour qu'un morphisme $f:X\to Y$ soit séparé, il faut que pour tout ouvert $U$
sur lequel $Y$ induit un préschéma séparé, le préschéma induit $f^{-1}(U)$
soit séparé, et il suffit qu'il en soit ainsi pour tout ouvert affine
$U\subset Y$.
\end{corollary}

La nécessité de la condition résulte de
(\hyperref[I.5.5.4-fr]{5.5.4}) et
(\hyperref[I.5.5.1-fr]{5.5.1, (ii)})~; la suffisance résulte de
(\hyperref[I.5.5.4-fr]{5.5.4}) et
(\hyperref[I.5.5.8-fr]{5.5.8}), compte tenu de l'existence de recouvrements
ouverts affines de $Y$.

En particulier, si $X$ et $Y$ sont des schémas affines, \emph{tout} morphisme
$X\to Y$ est séparé.

\begin{proposition}[5.5.10]
\label{I.5.5.10-fr}
Soient $Y$ un schéma, $f:X\to Y$ un morphisme. Pour tout ouvert affine $U$ de
$X$ et tout ouvert affine $V$ de $Y$, $U\cap f^{-1}(V)$ est affine.
\end{proposition}

Soient $p_1$, $p_2$ les projections de $X\times_{\mathbf{Z}}Y$~; le
sous-espace $U\cap f^{-1}(V)$ est l'image par $p_1$ de
$\Gamma_f(X)\cap p_1^{-1}(U)\cap p_2^{-1}(V)$. Or,
$p_1^{-1}(U)\cap p_2^{-1}(V)$ s'identifie à l'espace sous-jacent au

